
\begin{lemma}[Growth certificates for box QPs]\label{lem:growthcert}
Let $F=\frac12x^\T Hx+b^\T x+c$ be a quadratic on a continuous box $X=\bar X$, $x^*\in X$, $\zeta=\nabla F(x^*)$,
$S=\{i:\ell_i<x^*_i<u_i\}$ and $A=\{i:x^*_i\in\{\ell_i,u_i\}\}$.
\begin{enumerate}
\item[(a)] Suppose $\zeta_i=0$ for $i\in S$, $\zeta_i\ge0$ if $x_i^*=\ell_i$, and
$\zeta_i\le0$ if $x^*_i=u_i$. Let $M=\operatorname{diag}(\mu)$ with
$\mu_i=|\zeta_i|/(u_i-\ell_i)$ for $i\in A$ and $\mu_i=0$ for $i\in S$. Then
\begin{equation}\label{eq:minorant}
F(x)-F(x^*)\ge\tfrac12(x-x^*)^{\T}(H+2M)(x-x^*)\qquad(x\in X).
\end{equation}
Consequently, if $H+2M\succeq2\gamma I$ for some $\gamma>0$, then $x^*$ is the
unique minimizer and $F(x)-\OPT\ge\gamma\norm{x-x^*}^2$ on $X$.
\item[(b)] Conversely, if $x^*$ is a minimizer and
$F(x)-\OPT\ge g\norm{x-x^*}^2$ on $X$ with $g>0$, then $\zeta_S=0$ and
$H_{SS}\succeq2gI$. In particular $g\le\frac12v^{\T}Hv/\norm v^2$ for
every nonzero $v$ supported on $S$; if $S\ne\emptyset$ then every valid
curvature bound satisfies $L\ge2g$, so $\kappa\ge2$; and if $S=\{1,\dots,n\}$
then $H\succeq2gI$, so $F$ is strongly convex.
\end{enumerate}
\end{lemma}

\begin{proof}
(a) Let $d=x-x^*$. Since $F$ is quadratic,
$F(x)-F(x^*)=\ell^{\T}d+\frac12d^{\T}Hd$ exactly. For $i\in S$,
$\zeta_id_i=0$. For $i\in A$ with $x_i^*=\ell_i$ we have $0\le d_i\le u_i-\ell_i$
and $\zeta_i\ge0$, so $\zeta_id_i=|\zeta_i|d_i\ge|\zeta_i|d_i^2/(u_i-\ell_i)=\mu_id_i^2$;
the case $x_i^*=u_i$ is symmetric. Summing gives~\eqref{eq:minorant}. If
$H+2M\succeq2\gamma I$, the right side is at least $\gamma\norm d^2$, which is
positive for $d\ne0$.

(b) For $i\in S$ both directions $\pm e_i$ are feasible from $x^*$, so
minimality gives $\zeta_i=0$. Let $v\ne0$ be supported on $S$. For all
sufficiently small $|t|$, $x^*+tv\in X$, and
$F(x^*+tv)-\OPT=\frac{t^2}2v^{\T}Hv\ge gt^2\norm v^2$. Hence
$H_{SS}\succeq2gI$. Taking $v=e_i$ with $i\in S$ gives $H_{ii}\ge2g$; every
valid curvature bound satisfies $L\ge H_{ii}$, because
$t\mapsto F(\dots,t,\dots)-Lt^2/2$ has second derivative $H_{ii}-L$. If
$S$ is everything, $H=H_{SS}\succeq2gI$.
\end{proof}

\begin{remark}[Where nonconvexity can live]\label{rem:nonconvex}
Part~(b) says that for a continuous box QP, quadratic growth at an interior
minimizer forces strong convexity, and in general the Hessian block of the
free coordinates must be positive definite. A nonconvex continuous instance
covered by the growth theorem therefore has its negative curvature in
directions that involve coordinates at bounds, where first-order terms
$\zeta_id_i$ provide the growth, or in integer coordinates. Part~(a) is a
sufficient certificate of exactly this type. For an integer coordinate one
may also use $\zeta_id_i\ge-|\zeta_i|d_i^2$, valid because $|d_i|\le d_i^2$
for $d_i\in\Z$; this replaces $\mu_i$ by $-|\zeta_i|$. The planted instances of
Section~\ref{sec:computation} are built from part~(a): $H_{SS}$ has a
prescribed smallest eigenvalue, the entries touching $A$ make $H$ indefinite,
and $\mu$ is chosen so that $H+2M\succeq2\gamma I$ is verified exactly; then
$\gamma\le g\le\frac12v^{\T}H_{SS}v/\norm v^2$ for any rational $v$,
which brackets $\kappa$. Such instances are benign in one sense: once $x^*$
is known, \eqref{eq:minorant} certifies it directly.
\end{remark}

